\documentclass[10pt,a4paper]{amsart}
\usepackage[margin=19mm]{geometry}
\usepackage{amsmath,amssymb,mathtools}
\usepackage[scr=rsfs,cal=euler]{mathalfa}
\usepackage{xfrac}
\usepackage[unicode,hidelinks,bookmarks=false]{hyperref}
\newcommand{\ie}{i.e.\ }
\newcommand{\cf}{cf.\ }
\newcommand{\resp}{respectively\ }
\newcommand{\Subsection}{\subsection}
\providecommand{\nobreakdash}{\nobreak\hbox{-}}
\setlength{\emergencystretch}{2em}
\multlinegap0pt
\usepackage{tikz}
\newtheorem{theorem}{Theorem}
\title{Sequence b-colorings in graphs}
\author{Marko Jakovac, Michael S. Lang}
\date{Source: arXiv:2609.08484v1 (CC BY 4.0)}
\begin{document}
\maketitle

\noindent\emph{Source and licence.} Sequence b-colorings in graphs. Version arXiv:2609.08484v1 (CC BY 4.0), distributed by arXiv under the Creative Commons Attribution 4.0 International licence, http://creativecommons.org/licenses/by/4.0/, as recorded in the arXiv licence metadata for this exact version and shown on the abstract page https://arxiv.org/abs/2609.08484v1. Full licence notice: \texttt{licenses/arxiv-2609.08484-CC-BY-4.0.txt}.\par\smallskip

\noindent\textit{Editorial scope.} Counted unit: the article's theorem t:chi3 (source0.tex lines 768--771). The article's complete original proof of it is delivered with it (source0.tex lines 773--931, one \texttt{proof} environment of 159 lines). No mathematics is written, repaired or re-derived: the statement and the proof are the article's own lines, in the article's own notation, and no retained line is substituted or re-flowed.  Retained uncounted as the source-local context the proof needs: lines 85--85.  Nothing was omitted from any retained span, so the package carries no omission record.  No retained line contains a citation, so the wrapper carries no bibliography and no bibliography entry.  No retained line contains a figure environment, an image-inclusion command or drawing code, so no image is delivered and the package declares no asset dependency.  Editorial formatting only, in addition: an AMS \texttt{amsart} preamble that restates the article's own declarations and macros used by the retained lines, namely the environments \texttt{theorem} on separate counters, together with the standard packages those lines need.  The retained environments print in this document as Theorem 1 in order of appearance, whereas the article prints the same items with the numbering of its own full document.  The complete native source of the article as distributed in the arXiv e-print for this exact version is bundled unchanged as \texttt{sources/209717/source0.tex}.
\section{Definitions and initial results}\label{s:defs}

\begin{theorem}\label{t:chi3}
Let $G$ be a connected graph with $\chi(G)=3$. 
Then $G$ realizes the sequence $(2,1^\infty)$ if and only if $G$ is neither $C_5$ nor $T_{r,s,t}$ for any nonnegative integers $r$, $s$, and $t$. 
\end{theorem}

\begin{proof}
We first prove the forward implication by contraposition. Suppose that $G$ is either $C_5$ or a graph $T_{r,s,t}$ for some nonnegative integers $r$, $s$, and $t$. As shown before the theorem, both $C_5$ and every graph $T_{r,s,t}$ have chromatic number and b-chromatic number equal to $3$, but do not realize the sequence $(2,1,1)$. Since this is the only possible sequence of the form $(2,1^k)$ that they could realize, neither $C_5$ nor any graph $T_{r,s,t}$ realizes $(2,1^\infty)$.

For the converse, suppose that $G$ is a connected graph with $\chi(G)=3$ which is neither $C_5$ nor any graph $T_{r,s,t}$. We show that $G$ realizes $(2,1,1)$, and hence also $(2,1^\infty)$.

Since $\chi(G)=3$, there exists a proper coloring of $G$ using three colors.
As discussed in Section~\ref{s:defs}, this must be a b-coloring.
Consider such a coloring.
If one of the three color classes already contains at least two CDVs, then this coloring realizes $(2,1,1)$, and hence also $(2,1^\infty)$. Therefore, we may assume that each color class contains exactly one CDV. Let $x$, $y$, and $z$ denote the CDVs of colors $1$, $2$, and $3$, respectively. The subgraph induced by these three vertices is one of $K_3$, $P_3$, $\overline{P_3}$, or $\overline{K_3}$. We consider these four cases separately.

\smallskip

Suppose first that the vertices $x$, $y$, and $z$ induce $K_3$. If every neighbor of $x$, $y$, and $z$ outside the induced $K_3$ is a leaf, then, since $G$ is connected, $G$ is obtained from the induced $K_3$ by attaching leaves to its vertices. Hence $G$ is $T_{r,s,t}$ for some nonnegative integers $r$, $s$, and $t$, contrary to our assumption. 
Therefore, at least one of the vertices $x$, $y$, and $z$ has a non-leaf neighbor outside the induced $K_3$. 
Permuting colors if necessary,
let $u$ be a non-leaf neighbor of $z$ such that $u$ has color $1$. 
Since $u$ is not a leaf, it has a neighbor $v\neq z$.
As $u$ is not a CDV, it has no neighbor of color $2$, so
$v$ must
have color $3$.
This means that $v\not\in\{x,y\}$.
Since $v$ is not the 
unique
CDV of color $3$, it has no neighbor of color $2$.
We can thus recolor $v$ with color $2$ while preserving a proper coloring.
This keeps $x$, $y$ and $z$ as CDVs and makes $u$ a second CDV of color $1$.
The modified coloring thus realizes $(2,1,1)$ and hence also $(2,1^\infty)$.

\smallskip
Suppose next that $x$, $y$ and $z$ induce $P_3$.
Permuting colors if necessary, we may assume that $y$ is adjacent to $x$ and $z$.
(Recall that $x$, $y$ and $z$ have colors 1, 2 and 3, respectively.)
Since $x$ is a CDV, it is adjacent to a vertex $u$ that has color 3.
Similarly, $z$ is adjacent to a vertex $v$ that has color 1.
Let $A$ denote $\{u,x,y,z,v\}$ and set $A'=V(G)\setminus A$.
We proceed based on whether or not $u$ is adjacent to $v$.

Suppose that $u$ is adjacent to $v$, so $A$ induces a 5-cycle.
Since $G$ is assumed not to be $C_5$ itself, $A'$ is nonempty.
We distinguish three cases:
a1) at least one vertex in $A$ has a neighbor in $A'$ that is adjacent to another vertex in $A'$;
a2) no vertex in $A'$ has neighbors in both $A$ and $A'$ but at least one pair of vertices in $A$ have a common neighbor in $A'$;
a3) every vertex in $A'$ is a leaf.
In each case, we will show that $G$ can be recolored to realize $(2,1,1)$.

a1)
Suppose vertex $w_1\in A$ is adjacent to $w_2\in A'$, which is adjacent to $w_3\in A'$.
Permuting colors if necessary, we may assume that $w_1$ has color 1 and $w_2$ has color 2.
Since $w_2$ is not a CDV, $w_3$ must have color 1.
Since $w_3$ is not a CDV, it cannot have any neighbors of color 3.
We can thus recolor $w_3$ with color 3 while preserving properness of the coloring or changing the CDV status of $x$, $y$ or $z$.
Moreover, this recoloring turns $w_2$ into an extra CDV, as required.

a2)
Such a pair of vertices must have the same color, since otherwise the common neighbor would be a CDV.
Swapping colors 1 and 3 if necessary, we may assume that $x$ and $v$ are both adjacent to $w\in A'$.
We will recolor $G$.
Give vertices $u$, $x$, $y$, $z$, $v$ and $w$ the colors 1, 2, 3, 1, 3 and 1, respectively.
Any remaining vertices can be colored greedily since they are either leaves or perhaps adjacent only to $u$ and $z$.
Now we have a proper coloring where $u$ and $w$ are CDVs of color 1 and where $x$ and $y$ are CDVs of colors 2 and 3, respectively.
This new coloring thus realizes $(2,1,1)$.

a3)
Suppose $w\in A$ is adjacent to the leaf $w'\in A'$.
We will recolor $G$.
Starting at $w$ and going around the cycle, use the colors $1,2,3,1,2$.
This produces three CDVs in $A\setminus\{w\}$, one of each color.
Now give $w'$ color 3.
Any remaining vertices can be colored without causing conflicts.
Moreover, $w$ is a second CDV of color 1, as required.

Now suppose that $u$ is not adjacent to $v$, so $A$ induces a path.
We again distinguish three cases:
b1) at least one of $u$ or $v$ has a neighbor in $A'$;
b2) vertices $u$ and $v$ are leaves, but another vertex in $A$ has a non-leaf neighbor in $A'$;
b3) no vertex in $A$ has a non-leaf neighbor in $A'$.
We show that in cases b1 and b2, $G$ can be recolored to realize $(2,1,1)$ and that case b3 does not occur.

b1)
Suppose $u$ is adjacent to $w\in A'$.
Since $u$ is not a CDV, $w$ has color 1.
Because $w$ is not a CDV, it cannot have any neighbors of color 2.
We can thus recolor $w$ with color 2 while preserving the propriety of the coloring.
This turns $u$ into an extra CDV, as required.
The case in which $v$ has a neighbor in $A'$ is symmetric.

b2)
Let $w_1\in\{x,y,z\}$ be adjacent to $w_2\in A'$, which is itself adjacent to $w_3\neq w_1$.
Permuting colors as necessary, let $w_1$ have color 1 and $w_2$ have color 2.
Since $w_2$ is not a CDV, $w_3$ must have color 1.
This means that $w_3\in A'$, since $u$ and $v$ are leaves and $w_1$ is the only other vertex in $A$ that has color 1.
Because $w_3$ is not a CDV, it cannot have any neighbors of color 3.
Recolor $w_3$ with color 3.
This preserves the propriety of the coloring.
It also keeps $x$, $y$ and $z$ as CDVs.
Moreover, it turns $w_2$ into an extra CDV, as required.

b3)
In this case, $G$ is bipartite, contradicting the assumption $\chi(G)=3$.

\smallskip

Suppose now that $x$, $y$, and $z$ induce $\overline{P_3}$.
Recall that these vertices are so named so that they have colors 1, 2 and 3, respectively.
Permuting colors if necessary, we may assume that $y$ and $z$ are adjacent, while $x$ is adjacent to neither of them.
Let $X_2$ and $X_3$ denote the sets of neighbors of $x$ with colors $2$ and $3$, respectively. Since $x$ is a CDV, both sets are nonempty. Let $Y_1$ denote the set of neighbors of $y$ with color $1$, and let $Z_1$ denote the set of neighbors of $z$ with color $1$. These sets are also nonempty because $y$ and $z$ are CDVs. 

We may assume that every neighbor of $y$ different from $z$ has color $1$ and is thus in $Y_1$.
Indeed, suppose that $v\neq z$ is a neighbor of $y$ with color $3$.
Since $v$ is not a CDV and is adjacent to $y$ of color $2$, $v$ has no neighbor of color $1$.
We may thus give $v$ color 1 without disturbing the propriety of the coloring or the CDV status of $x$, $y$ and $z$.
By symmetry, we may also assume that every neighbor of $z$ different from $y$ belongs to $Z_1$.

We may also assume that we have accounted for all of the vertices, so $V(G)=\{x,y,z\}\cup X_2\cup X_3\cup Y_1\cup Z_1$.
For example, suppose that a vertex $v\in X_2$ has a neighbor $u\notin \{x,y,z\}\cup X_2\cup X_3\cup Y_1\cup Z_1$. 
Since $v$ (with color 2) is not a CDV and is adjacent to $x$ (of color $1$), $u$ must have color $1$.
Moreover, $u$ is not a CDV and is adjacent to $v$ of color $2$, so it has no neighbor of color $3$. We may therefore recolor $u$ with color $3$, making $v$ a second CDV of color $2$. 
This modified coloring realizes $(2,1,1)$, and hence also $(2,1^\infty)$.
By symmetry, the same argument applies to every vertex in $X_3$, $Y_1$, and $Z_1$.

Let us consider the edges among the four sets $X_2$, $X_3$, $Y_1$ and $Z_1$. 
There are no edges inside any one of them, since each set is monochromatic.
There are also no edges between $Y_1$ and $Z_1$, since all vertices in these two sets have color $1$. 
An edge between $X_2$ and $X_3$ would make both endpoints CDVs.
An edge between $X_2$ and $Z_1$ would make its endpoint in $Z_1$ a second CDV of color $1$, while an edge between $X_3$ and $Y_1$ would make its endpoint in $Y_1$ a second CDV of color $1$.
Therefore, the only possible edges among these sets are between $X_2$ and $Y_1$ and between $X_3$ and $Z_1$.
There must be at least one edge of the latter type.
Indeed, if no vertex in $X_3$ were adjacent to any vertex in $Z_1$, then we could properly recolor the vertices so those in $X_3\cup X_2\cup\{y\}\cup Z_1$ have color 1 and the rest have color 2.
But this would show that $G$ is bipartite, contradicting the assumption $\chi(G)\geq3$.

We now recolor some vertices.
Assign color 2 to every vertex in $Z_1$, color 1 to $z$ and color 3 to $y$.
By the above description of the possible edges of $G$, this coloring is proper. 
Choose an edge $uv$ with $u\in X_3$ and $v\in Z_1$.
Then $u$ is a CDV of color $3$, since it is adjacent to $x$ of color $1$ and to $v$ of color $2$, while $v$ is a CDV of color $2$, since it is adjacent to $z$ of color $1$ and to $u$ of color $3$. 
We have not recolored the neighbors of $x$, so it is still a CDV of color 1.
However, $z$ is now adjacent to $y$ of color 3 and $v$ of color 2, so it is a second CDV of color 1.
This new coloring realizes $(2,1,1)$, and hence also $(2,1^\infty)$.

\smallskip

Suppose finally that the vertices $x$, $y$, and $z$ induce $\overline{K_3}$. Recall that $x$, $y$, and $z$ have colors $1$, $2$, and $3$, respectively. Let $X_2$ and $X_3$ denote the sets of neighbors of $x$ with colors $2$ and $3$, respectively. Similarly, let $Y_1$ and $Y_3$ denote the sets of neighbors of $y$ with colors $1$ and $3$, and let $Z_1$ and $Z_2$ denote the sets of neighbors of $z$ with colors $1$ and $2$. Since $x$, $y$, and $z$ are CDVs, all six sets are nonempty. 
These sets are pairwise disjoint, since any vertex in two of them would be an extra CDV.

As in the previous case, we consider whether there are any more vertices.
Suppose, for example, that a vertex $v\in X_2$ has a neighbor
$u\notin\{x,y,z\}\cup X_2\cup X_3\cup Y_1\cup Y_3\cup Z_1\cup Z_2$.
Since $v$ is not a CDV and is adjacent to $x$ of color $1$, it has no neighbor of color $3$. Thus $u$ has color $1$. Moreover, $u$ is not a CDV and is adjacent to $v$ of color $2$, so it has no neighbor of color $3$. We may therefore recolor $u$ with color $3$, making $v$ a second CDV of color $2$. This modified coloring realizes $(2,1,1)$, and hence also $(2,1^\infty)$. The same argument applies to every vertex in the other five sets. Therefore, we may assume that
$V(G)=\{x,y,z\}\cup X_2\cup X_3\cup Y_1\cup Y_3\cup Z_1\cup Z_2$.

There are clearly no edges between vertices of the same color. 
An edge between $X_2$ and $Y_3$ would make its endpoint in $X_2$ a CDV of color $2$.
Analogous statements hold for most of the other pairs.
In the end, we find that the only possible edges among the six sets are between $X_2$ and $Y_1$, between $X_3$ and $Z_1$, and between $Y_3$ and $Z_2$. 
Since $G$ is connected, at least two of these three types of edge must be present.
By symmetry, assume that there is an edge between $X_2$ and $Y_1$ and an edge between $X_3$ and $Z_1$.

We now define a new coloring. Assign color $1$ to $x$, $z$, and every vertex of $Y_1\cup Y_3$. Assign color $2$ to every vertex of $X_2\cup Z_1\cup Z_2$. Finally, assign color $3$ to $y$ and every vertex of $X_3$. By the above description of the possible edges, this coloring is proper. Choose an edge $uv$ with $u\in X_2$ and $v\in Y_1$. Then $v$ is a CDV of color $1$, since it is adjacent to $u$ of color $2$ and to $y$ of color $3$. The vertex $x$ is also a CDV of color $1$, since it has neighbors in both $X_2$ and $X_3$. Similarly, choose an edge $wp$ with $w\in X_3$ and $p\in Z_1$. Then $w$ is a CDV of color $3$, since it is adjacent to $x$ of color $1$ and to $p$ of color $2$. The vertex $p$ is a CDV of color $2$, since it is adjacent to $z$ of color $1$ and to $w$ of color $3$. This new coloring realizes $(2,1,1)$, and hence also $(2,1^\infty)$.
\end{proof}

\end{document}
